% ============================================================
%  Abstract – Modulbeschreibung
% ============================================================
\chapter*{Abstract}
\addcontentsline{toc}{chapter}{Abstract}
\markboth{Abstract}{Abstract}

\begin{tabular}{@{}ll@{}}
  \textbf{Modul:}       & Systemtheorie 1 – Signale und Systeme (141170) \\[0.3em]
  \textbf{Dozent:}      & Prof.\,Dr.-Ing.\ Rainer Martin \\[0.3em]
  \textbf{Fakultät:}    & Elektrotechnik und Informationstechnik, Ruhr-Universität Bochum \\[0.3em]
  \textbf{Turnus:}      & jedes Sommersemester \quad|\quad 4\,SWS \quad|\quad 5\,CP \\[0.3em]
  \textbf{Vorkenntnisse:} & Komplexe Zahlen, Differential- und Integralrechnung, Reihen \\
\end{tabular}

\vspace{1.2em}

Die Vorlesung \textbf{Systemtheorie 1 – Signale und Systeme} vermittelt die mathematischen
Grundlagen zur Beschreibung, Analyse und Verarbeitung von Signalen sowie zum systematischen
Umgang mit linearen und nichtlinearen Systemen. Sie bildet das Fundament für weiterführende
Veranstaltungen in der Nachrichten-, Mess- und Regelungstechnik.

\vspace{0.8em}

\section*{Modulinhalt}

Das Modul gliedert sich in drei thematische Blöcke:

\subsection*{1\quad Signale und Systeme}

Der erste und umfangreichste Teil behandelt die mathematische Beschreibung von Signalen
im Zeitbereich. Ausgangspunkt sind grundlegende Kenngrößen und Eigenschaften von Signalen
– Energie, Leistung, Periodizität und Symmetrie – sowie elementare Operationen wie
Verschiebung, Spiegelung und Skalierung. Auf dieser Basis werden Methoden der
\textbf{Signalsynthese und -analyse} eingeführt, insbesondere die Zerlegung periodischer
Signale in Fourierreihen. Ein weiterer Schwerpunkt liegt auf der
\textbf{Analog-Digital- und Digital-Analog-Umsetzung} (Abtasttheorem nach
Shannon–Nyquist, Quantisierung). Den Abschluss bildet die Klassifikation von Systemen:
lineare vs.\ nichtlineare, zeitinvariante vs.\ zeitvariante sowie kausale Systeme und
deren Beschreibung durch Impulsantwort und Faltungsintegral.

\subsection*{2\quad Einführung in die Wahrscheinlichkeitsrechnung}

Der zweite Block legt die stochastischen Grundlagen, die für die spätere Behandlung
von Rauschen und Zufallssignalen unverzichtbar sind. Behandelt werden Axiome und
Rechenregeln der Wahrscheinlichkeit, bedingte Wahrscheinlichkeit und Bayes'sches
Theorem sowie mehrstufige Zufallsexperimente. Im Mittelpunkt stehen
\textbf{diskrete und kontinuierliche Zufallsvariablen} mit ihren Verteilungsfunktionen,
Dichtefunktionen und charakteristischen Kenngrößen (Erwartungswert, Varianz, Momente).

\subsection*{3\quad Grundbegriffe der Informationstheorie}

Der dritte Block führt in die von Claude Shannon begründete Informationstheorie ein.
Zentrale Fragen sind: Was ist Information? Wie lässt sie sich messen? Wie effizient
kann eine Quelle komprimiert werden? Als mathematisches Werkzeug wird der
\textbf{Entropiebegriff} (Shannon-Entropie, bedingte Entropie, Transinformation)
eingeführt und auf praktische Anwendungen in der Quellen- und Kanalcodierung übertragen.

\vspace{0.8em}

\section*{Lernziele}

\begin{itemize}
  \item Mathematische Beschreibung von Signalen und Systemen im Zeitbereich und
        Kenntnis ihrer wesentlichen Merkmale.
  \item Sicherer Umgang mit diskreten und kontinuierlichen Zufallsvariablen sowie
        deren Kenngrössen.
  \item Verständnis der grundlegenden Konzepte der Informationstheorie und Fähigkeit
        zur Anwendung des Entropiebegriffs.
  \item Vernetzung von Signaltheorie, Stochastik und Informationstheorie als Basis
        für weiterführende Gebiete der Elektrotechnik und Informationstechnik.
\end{itemize}
